% problem-000182 4 尿碘Sandell反应检测
\section*{4. 尿碘Sandell反应检测}
斜⽅辉⽯[ $\left( \mathrm { M g , F e ) _ { 2 } S i _ { 2 } O _ { 6 } } \right]$ 是地壳和上地幔的主要组成矿物之⼀，在其晶体结构中包含有两种不同的硅氧四⾯体 $( \mathrm { S i O _ { 3 } } ) ^ { 2 }$ −链，分别称作A链和B链，基于这两种链的排布⽽形成了两种结构有差异的⼋⾯体空隙M1和M2，⼆者⽐例相同，如下图所⽰，Mg和 $\mathrm { F e ^ { 2 } }$ 便分布在这些⼋⾯体空隙中。

（图：）
(a)

（图：）
(b)

由于M1和M2的空隙性质不同，导致 $\mathrm { M g ^ { 2 + } }$ 和 $\mathrm { F e ^ { 2 + } }$ 对其占据的选择性不同， $\mathrm { F e ^ { 2 + } }$ 倾向于占据M2位置。⼀定条件下，两种离⼦可以发⽣不同位置的交换反应：

$
\mathrm{Fe} ^ {2 +} (\mathrm{M} 1) + \mathrm{Mg} ^ {2 +} (\mathrm{M} 2) \rightleftharpoons \mathrm{Fe} ^ {2 +} (\mathrm{M} 2) + \mathrm{Mg} ^ {2 +} (\mathrm{M} 1)
$

为简便便起⻅， $\mathrm { F e ^ { 2 + } ( M l ) }$ $\mathrm { F e } ^ { 2 + } ( \mathrm { M } 2 )$ $\mathrm { M g ^ { 2 + } ( M l ) }$ 、 $\mathrm { M g ^ { 2 + } ( M 2 ) }$ 分别写作 ${ \mathrm { F e } } ( 1 )$ ${ \mathrm { F e } } ( 2 )$ 、Mg(1)、 $\mathrm { M g } ( 2 )$ 。上述反应达平衡时，分配系数 $K _ { D }$ ：

$
K _ {D} = \frac {\chi_ {\mathrm{Fe(2)}} \chi_ {\mathrm{Mg(1)}}}{\chi_ {\mathrm{Fe(1)}} \chi_ {\mathrm{Mg(2)}}}
$

其中， $\chi$ 为各离⼦占据相应位置的摩尔分数，例如： $\chi _ { \mathrm { F e } ( 2 ) }$ 为 $\mathrm { F e ^ { 2 + } }$ 离⼦占据M(2)位置的摩尔分数，其他同理。选择某⼀矿物样品，在873K下进⾏处理，利⽤X射线衍射结合穆斯堡尔谱监测反应进⾏过程中上述物种占据不同位置情况随时间的变化，数据列⼊下表中。

\textit{（原卷此处含表格/仪器试剂表，详见 PDF 或 MD 源文件。）}

\subsection*{4-1}
计算分配系数 $K _ { D ^ { \varsigma } }$ 。（提⽰：根据表中数据，合理判断并选择数据）。

\subsection*{4-2}
上述反应的正、逆反应的速率常数分别为 $| k _ { 1 }$ 和� ，假设正、逆反应速率表达形式均与基元反应类似，正、逆反应速率分别与占据相应位置的各离⼦的摩尔分数�成正⽐。写出 $K _ { D }$ 与 $k _ { 1 }$ 和� 的关系 $\updownarrow \updownarrow _ { \updownarrow }$

\subsection*{4-3}
利⽤上表中起始阶段的数据（要求：采⽤编号1\~4的数据进⾏处理），计算 $k _ { 1 }$ 和� 的值。

提⽰1：可将 $\mathrm { M g ^ { 2 } }$ +的摩尔分数近似视为常数，取0.9780；

提⽰2：常微分⽅程的解，对于�(�)有：

$
y ^ {\prime} = k y, \text {其解为:} \ln y = k x + c
$

$
y ^ {\prime} = k (a y + b), \text {其解为:} \frac {1}{a} \mathrm{ln} (a y + b) = k x + c
$
